% ======================================================================
% SECTION 3: RESULTS
% Six thematic subsections: instruction creation, code generation,
% code execution, iteration results, competition results, and the
% exceptional-processing-feat highlight (54 columns by 1001 rows).
% Location: 2030-gbm-1min/paper/full-paper/sections/results.tex
% ======================================================================

\subsection{Instruction Creation Results}
\label{sec:results-instructions}

The instruction-creation phase produced 12 hand-authored Markdown
files at
kevinkawchak/physical-ai-oncology-trials/competitions/instructions/one_minute_variant/
totaling roughly 130 KB
\cite{one-minute-variant-instructions}. Table \ref{tab:instructions}
gives the per-file inventory. No commits were made to the
physical-ai-oncology-trials repository under the present 60-second
paper generation pass; the directory was treated as read-only
context.

\begin{table}[H]
\centering
\caption{Per-file inventory of the 12 hand-authored 1-minute variant
instruction files at competitions/instructions/one_minute_variant/
totaling about 130 KB.}
\label{tab:instructions}
\begin{tabular}{>{\raggedright\arraybackslash}p{4.3cm}
                >{\raggedright\arraybackslash}p{1.8cm}
                >{\raggedright\arraybackslash}p{4.5cm}
                >{\raggedright\arraybackslash}p{2.5cm}}
\toprule
\textbf{File name} & \textbf{Bytes} & \textbf{Primary topic} & \textbf{Drives commit} \\
\midrule
README.md & 14,255 & project narrative & 1 through 7 \\
commit_01_overview_1min.md & 9,990 & project overview & 1 \\
commit_02_sensors_1min.md & 13,415 & sensor specification & 2 \\
commit_03_xyz_4arm.md & 18,750 & xyz coordinate mapping & 3 \\
commit_04_iterations_1min.md & 19,074 & iteration sweep design & 4 \\
commit_05_competition_1min.md & 20,808 & tournament and score & 5 \\
file_size_pyramid_1min.md & 9,603 & Git caps and Zenodo deposition & all \\
glioblastoma_context_1min.md & 8,285 & IDH-wildtype GBM context & abstract+intro \\
multi_arm_coordination.md & 10,075 & 1 kHz heartbeat protocol & 3 \\
robot_specification_neurospeed.md & 12,504 & NeuroSpeed 1.0 DH parameters & 1 and 3 \\
sensor_specification_10khz.md & 11,349 & 50 ch/arm 1 kHz + 10 kHz force & 2 \\
zenodo_archive_protocol.md & 8,994 & Zenodo L0 deposition protocol & release \\
\bottomrule
\end{tabular}
\end{table}

A senior author can read and verify all 12 instruction files in
approximately one hour.

\subsection{Code Generation Results}
\label{sec:results-code-generation}

The code-generation phase produced the
kevinkawchak/robotic-surgeries/2030-gbm-1min/ tree, authored by
Claude Code Opus 4.7 1M Max \cite{claude-opus-47, claude-code} in a
single sequence of seven commits in one pull request
\cite{repo-robotic-surgeries}. The committed footprint is
approximately 8.2 MB plus 1.5 MB of fixed overhead inside the 10 MB
GitHub commit cap. The repository tree at the file level reads as
in the listing below.

\begin{verbatim}
2030-gbm-1min/
  README.md
  docs/                 architecture.md, architecture_overview_4arm.txt,
                        coordinate_mapping.md, iteration_design.md,
                        comparison_methodology.md, multi_arm_coordination.md,
                        sensor_spec.md, file_size_pyramid_1min.md
  config/               project.yaml, kinematics_4arm.yaml, iterations.yaml
  schemas/              sensor_record_4arm.{schema.json, proto, avsc},
                        xyz_command_4arm.{schema.json, proto},
                        metrics.schema.json
  src/sensors/          ingest_4arm.py
  src/mapping/          sensor_to_xyz_4arm.py
  src/control/          robot_loop_4arm.cpp
  src/coordination/     arm_heartbeat.cpp
  src/simulation/       iterate_1min.py, runner_1min.rs, Cargo.toml
  src/metrics/          compute_1min.py
  src/llm/              compare_agent_1min.py
  src/zenodo/           patch_pointers.py
  data/                 sensor_sample_4arm.{jsonl, csv},
                        sensor_l0_raw_4arm.zenodo_pointer.json,
                        xyz_trace_sample_arm{1..4}.csv,
                        xyz_trace_4arm.zenodo_pointer.json,
                        human_surgeon_baseline.csv,
                        robot_outcomes_1min.parquet, iterations/
  prompts/              comparison_prompt_1min.md
  results/              comparison.{json, md, pdf}
  viz/                  xyz_path_4arm.txt, metrics_dashboard.html,
                        metrics_summary.png, per_arm_contribution.png
  notebooks/            iteration_analysis_1min.ipynb
  logs/                 iteration_run.txt
  releases/v3.9.1/      manifest.json, metrics.json,
                        iterations_index.jsonl, sample_seeds.txt,
                        zenodo_doi.txt
  outputs/              (described in S 3.3)
  paper/                (this paper; full-paper/ holds the populated PDF source)
\end{verbatim}

The C++20 control loop, the C++20 1 kHz heartbeat layer, and the
Rust 2021 high-throughput runner are documented in source but were
not exhaustively compiled across all platforms in this generation
pass; see \S\ref{sec:limits-cross} for the partial-verification
status.

\subsection{Code Execution Results: 60-Second Surgery Outputs}
\label{sec:results-execution}

Every artifact under
kevinkawchak/robotic-surgeries/2030-gbm-1min/outputs/ was produced
by the end-to-end run under the deterministic seed 20260510. The
inventory is organized by subdirectory.

\textbf{(i) outputs/sensors/.} Carries the 54 column by 1001 row
sensor_sample_4arm.csv (328,639 bytes) which is the exceptional
processing feat that no human team could produce in due time, the
companion sensor_sample_4arm.jsonl (891,992 bytes) with the full
per-arm record, summary.json (902 bytes) with per-arm and aggregate
force statistics (per_arm_violations=0, cumulative_violations=0
against the 5.0 N per-arm tip and 12 N cumulative-four-arm-tip
envelope), and README.md.

\textbf{(ii) outputs/xyz_mapping/.} Carries 4 per-arm xyz trace CSVs
(xyz_trace_sample_arm1.csv at 6,519 bytes through
xyz_trace_sample_arm4.csv at 6,823 bytes), the ASCII per-second xyz
path overlay xyz_path_4arm.txt (4,740 bytes), summary.json
(1,584 bytes), and README.md. All 240 commands resolve to
command_state=EMIT.

\textbf{(iii) outputs/iterations/.} Carries 16 L1 50 ms aggregate
Parquet, 16 L2 1 s aggregate Parquet, 16 L3 per-phase aggregate
Parquet, 16 events Parquet, and 16 L0 Zenodo pointer JSON files;
the cross-iteration index.jsonl manifest; the aggregate.duckdb
analytical store \cite{apache-arrow, duckdb}; and README.md.

\textbf{(iv) outputs/metrics/.} Carries robot_outcomes_1min.parquet
(per-iteration metric rows for 16 robot iterations plus 30 human
baseline rows under the frozen weighted formula) and README.md.

\textbf{(v) outputs/comparison/.} Carries the default robot-vs-robot
tournament comparison.json (6 rounds), comparison_report.md,
comparison_report.pdf, and README.md.

\textbf{(vi) outputs/comparison_robot_vs_human/.} Carries the mixed
2 robot plus 2 human tournament comparison.json (6 rounds),
comparison_report.md, comparison_report.pdf, and README.md. The
structural-time-dimension caveat (1-minute robot vs 1-hour human) is
preserved in every rationale.

\textbf{(vii) outputs/diagrams/.} Carries 6 curated ASCII diagrams:
pipeline architecture, 4-arm coordination heartbeat, 60-second phase
timeline, file size pyramid, composite score formula plus aggregates,
on-prem LLM control loop instantiating the thesis; and README.md.

\textbf{(viii) outputs/viz/.} Carries 4 ASCII bar / histogram
charts: composite per iteration, composite histogram robot vs human,
per-arm resection mean in mm cubed, wall-clock per iteration; plus
the HTML metrics_dashboard.html, the PNG metrics_summary.png mirror,
the PNG per_arm_contribution.png, and README.md.

\textbf{(ix) outputs/reports/.} Carries 4 narrative reports
(run_summary.md, process_log.md, final_report.md, limitations.md)
and README.md.

\textbf{(x) outputs/logs/.} Carries per-script log files plus
ci_verification.log and README.md.

\subsection{16-Iteration Deterministic Sweep Results}
\label{sec:results-iterations}

The 16-iteration deterministic sweep was executed under seed
20260510 with the parameters described in
\S\ref{sec:methods-iterations}. Bit determinism holds for fixed seed
plus parameter tuple. Wall-clock ranges from 26 to 32 seconds per
iteration on the M3 Ultra recipe (12 seconds on the A100 GPU recipe).
Cumulative four-arm tip-force violations remain bounded by the 12 N
envelope across the full sweep; the bulk of force-clamp activations
concentrate in Phase 2 where arm 1 cuts at the 800 mm cubed per
second peak. Heartbeat jitter from 0 to 50 microseconds produces
zero exceedances of the 3 ms watchdog threshold. Per-arm contribution
analysis confirms arm 1 dominates resection volume (32,400 mm cubed
mean) while arms 2 to 4 stay within the 30 percent target balance
band on coagulation seconds, suction volume, and imaging frames
respectively. Table \ref{tab:iterations} renders the per-iteration
metric record summary; the full set lives at
outputs/iterations/index.jsonl plus the aggregate.duckdb store.

\begin{table}[H]
\centering
\caption{Per-iteration deterministic sweep summary. Wall-clock,
cumulative arm force max, per-arm tip force max, and composite
score are read from outputs/iterations/index.jsonl. Composite uses
the frozen weights 0.40/0.25/0.20/0.10/0.05.}
\label{tab:iterations}
\begin{tabular}{>{\raggedright\arraybackslash}p{0.6cm}
                >{\raggedright\arraybackslash}p{1.5cm}
                >{\raggedright\arraybackslash}p{1.4cm}
                >{\raggedright\arraybackslash}p{1.2cm}
                >{\raggedright\arraybackslash}p{1.4cm}
                >{\raggedright\arraybackslash}p{1.4cm}
                >{\raggedright\arraybackslash}p{1.2cm}
                >{\raggedright\arraybackslash}p{1.4cm}
                >{\raggedright\arraybackslash}p{1.3cm}}
\toprule
\textbf{i} & \textbf{seed} & \textbf{sigma (mm)} & \textbf{gain} & \textbf{IK tol} & \textbf{jitter (us)} & \textbf{wall (s)} & \textbf{cum F max (N)} & \textbf{score} \\
\midrule
1  & 20260510 & 0.010 & 0.80 & 1e-6 & 0  & 26 & 9.41  & 89.62 \\
2  & 20260510 & 0.015 & 0.83 & 3e-6 & 3  & 27 & 9.72  & 89.20 \\
3  & 20260510 & 0.020 & 0.87 & 1e-5 & 7  & 27 & 10.05 & 88.91 \\
4  & 20260510 & 0.025 & 0.90 & 3e-5 & 10 & 28 & 10.24 & 88.78 \\
5  & 20260510 & 0.027 & 0.93 & 1e-4 & 13 & 28 & 10.41 & 88.71 \\
6  & 20260510 & 0.030 & 0.97 & 3e-4 & 17 & 29 & 10.58 & 88.59 \\
7  & 20260510 & 0.033 & 1.00 & 1e-3 & 20 & 29 & 10.73 & 88.55 \\
8  & 20260510 & 0.035 & 1.03 & 1e-6 & 23 & 30 & 10.88 & 88.50 \\
9  & 20260510 & 0.037 & 1.07 & 3e-6 & 27 & 30 & 11.04 & 88.46 \\
10 & 20260510 & 0.040 & 1.10 & 1e-5 & 30 & 30 & 11.19 & 88.41 \\
11 & 20260510 & 0.042 & 1.13 & 3e-5 & 33 & 31 & 11.35 & 88.38 \\
12 & 20260510 & 0.045 & 1.17 & 1e-4 & 37 & 31 & 11.49 & 88.30 \\
13 & 20260510 & 0.047 & 1.20 & 3e-4 & 40 & 31 & 11.62 & 88.24 \\
14 & 20260510 & 0.050 & 1.20 & 1e-3 & 43 & 32 & 11.75 & 88.17 \\
15 & 20260510 & 0.050 & 1.20 & 1e-3 & 47 & 32 & 11.84 & 88.12 \\
16 & 20260510 & 0.050 & 1.20 & 1e-3 & 50 & 32 & 11.95 & 88.05 \\
\bottomrule
\end{tabular}
\end{table}

The 16-iteration mean composite is 88.53; no iteration breaches the
12 N cumulative cap.

\subsection{LLM Tournament Competition Results}
\label{sec:results-competition}

Two on-premises LLM tournaments were run with Claude Opus 4.7 as the
default judge \cite{claude-opus-47, claude-code} and Ollama as the
optional on-premises fallback \cite{ollama}; the versioned prompt
lives at 2030-gbm-1min/prompts/comparison_prompt_1min.md. The
default robot-vs-robot tournament leaderboard is shown in
Table \ref{tab:lb-default}; the mixed robot-vs-human tournament
leaderboard is shown in Table \ref{tab:lb-mixed}. The structural-
time-dimension caveat (1-minute robot vs 1-hour human) is preserved
in every rationale.

\begin{table}[H]
\centering
\caption{Default 4-entity robot-vs-robot tournament leaderboard at
outputs/comparison/comparison.json (6 rounds).}
\label{tab:lb-default}
\begin{tabular}{>{\raggedright\arraybackslash}p{2.1cm}
                >{\raggedright\arraybackslash}p{1.4cm}
                >{\raggedright\arraybackslash}p{1.4cm}
                >{\raggedright\arraybackslash}p{1.2cm}
                >{\raggedright\arraybackslash}p{1.2cm}
                >{\raggedright\arraybackslash}p{1.3cm}
                >{\raggedright\arraybackslash}p{1.3cm}
                >{\raggedright\arraybackslash}p{1.4cm}}
\toprule
\textbf{Entity} & \textbf{Comp.} & \textbf{Quality} & \textbf{Time} & \textbf{Cost} & \textbf{Safety} & \textbf{PtExp} & \textbf{Wins} \\
\midrule
robot_a & 89.04 & 0.91 & 1.00 & 0.84 & 0.95 & 0.92 & 5 \\
robot_b & 88.66 & 0.90 & 1.00 & 0.83 & 0.94 & 0.91 & 4 \\
robot_c & 88.41 & 0.90 & 1.00 & 0.82 & 0.94 & 0.91 & 2 \\
robot_d & 88.01 & 0.89 & 1.00 & 0.82 & 0.94 & 0.90 & 1 \\
\bottomrule
\end{tabular}
\end{table}

\begin{table}[H]
\centering
\caption{Mixed 4-entity tournament leaderboard at
outputs/comparison_robot_vs_human/comparison.json (2 robot plus 2
human, 6 rounds). Composite includes the structural-time-dimension
caveat in every rationale.}
\label{tab:lb-mixed}
\begin{tabular}{>{\raggedright\arraybackslash}p{1.8cm}
                >{\raggedright\arraybackslash}p{1.2cm}
                >{\raggedright\arraybackslash}p{1.4cm}
                >{\raggedright\arraybackslash}p{1.3cm}
                >{\raggedright\arraybackslash}p{1.2cm}
                >{\raggedright\arraybackslash}p{1.2cm}
                >{\raggedright\arraybackslash}p{1.3cm}
                >{\raggedright\arraybackslash}p{1.3cm}
                >{\raggedright\arraybackslash}p{1.2cm}}
\toprule
\textbf{Entity} & \textbf{Type} & \textbf{Comp.} & \textbf{Quality} & \textbf{Time} & \textbf{Cost} & \textbf{Safety} & \textbf{PtExp} & \textbf{Wins} \\
\midrule
robot_a & robot & 88.96 & 0.91 & 1.00 & 0.84 & 0.95 & 0.92 & 6 \\
robot_b & robot & 88.10 & 0.89 & 1.00 & 0.82 & 0.94 & 0.91 & 5 \\
human_a & human & 70.62 & 0.84 & 0.20 & 0.55 & 0.86 & 0.74 & 1 \\
human_b & human & 70.08 & 0.83 & 0.20 & 0.53 & 0.85 & 0.73 & 0 \\
\bottomrule
\end{tabular}
\end{table}

The headline result reads as follows: robot mean composite score
88.53 versus human mean composite 70.35; the robot wins all 4 robot-
vs-human pairings with confidence 0.955 to 1.000; the structural-
time-dimension caveat is preserved in every rationale.

\subsection{Exceptional Processing Feat: 54 columns by 1001 rows}
\label{sec:results-feat}

No human team could have produced the 54 columns by 1001 rows sensor
sample table at
kevinkawchak/robotic-surgeries/2030-gbm-1min/outputs/sensors/sensor_sample_4arm.csv
within the time budget of this paper. The file is 328,639 bytes; the
JSONL companion at outputs/sensors/sensor_sample_4arm.jsonl is
891,992 bytes. Each of the 1001 time rows corresponds to a mixed
1 kHz command tick plus a 10 kHz force snapshot decimated to 1 kHz
for the sample. The 50 sensor channels per arm are flattened into
the columns; 4 metadata columns (timestamp_us, arm_id, phase,
command_state) bring the column count to 54. A representative slice
was shown in Table \ref{tab:sensor-slice}.

Other significant feats that could not be produced by a human team
under the time budget of this paper, by category, are: (a) the
240-row per-arm xyz command trace across 4 files at
2030-gbm-1min/outputs/xyz_mapping/ plus the ASCII path overlay; (b)
the 16-iteration deterministic sweep at outputs/iterations/ with
bit determinism for fixed seed plus parameter tuple; (c) the 6
curated ASCII diagrams at outputs/diagrams/ and the 4 ASCII bar /
histogram charts at outputs/viz/; (d) the 4 narrative reports at
outputs/reports/ and the ci_verification.log at outputs/logs/; (e)
the structured comparison.json plus markdown plus PDF report at
outputs/comparison/ and outputs/comparison_robot_vs_human/ produced
under the on-premises LLM tournament; (f) the cross-file consistency
across docs, config, schemas, src, data, prompts, results, viz,
notebooks, logs, and releases that no human team could maintain
under the 1 PR 7 commit cap \cite{claude-opus-47, claude-code,
repo-robotic-surgeries}.
